\section{Обзор литературы}

\paragraph{Диффузионные модели и механизмы guidance.}
Современные диффузионные модели являются стандартом для задач генерации изображений благодаря высокой фотореалистичности. Качество условной генерации в них во многом определяется механизмом \textit{guidance}~--- способом направленного изменения траектории обратного диффузионного процесса.

Ранние работы по \textit{classifier guidance} показали, что добавление градиента лог-вероятности класса к оценке шума позволяет управлять компромиссом между качеством и разнообразием генерации \cite{dhariwal2021diffusion}. Формально, если $\epsilon_\theta(x_t, t)$~--- оценка шума, а $\nabla_{x_t} \log p_\phi(y \mid x_t)$~--- градиент классификатора, то обновление принимает вид
\[
\tilde{\epsilon}_\theta = \epsilon_\theta + s \cdot \sigma_t \nabla_{x_t} \log p_\phi(y \mid x_t),
\]
где $s$~--- коэффициент guidance.

Позднее был предложен \textit{classifier-free guidance (CFG)}, позволяющий избежать обучения отдельного классификатора \cite{ho2022classifierfree}. В этом случае комбинируются условная и безусловная оценки шума:
\[
\epsilon_{\text{CFG}} = (1 + w) \epsilon_\theta(x_t, t, c) - w \epsilon_\theta(x_t, t, \emptyset)
\]
где $c$~--- текстовое условие, $\emptyset$~--- безусловный режим, $w$~--- коэффициент guidance.

Однако, несмотря на практическую эффективность, \textit{classifier-free guidance} в первую очередь усиливает уже присутствующую в модели условную компоненту, а не вводит новую семантическую информацию. По сути, коэффициент $w$ масштабирует разницу между условной и безусловной оценками шума, тем самым повышая выразительность признаков, коррелирующих с текстовым условием. Это приводит к улучшению семантической чёткости и согласованности с промптом, однако не гарантирует корректную реализацию сложных или редких композиционных структур. Если модель изначально не выучила определённую зависимость, например нетипичную пространственную конфигурацию или логически противоречивый атрибут, то увеличение $w$ лишь усиливает наиболее вероятные интерпретации, но не устраняет фундаментальные ограничения распределения. Таким образом, CFG частично решает задачу семантической точности за счёт усиления условного сигнала, но не обеспечивает полноценного контроля над сложными семантиками.

\paragraph{VLM-guidance и оптимизация латентов.}
Одним из направлений решения данной проблемы стало использование больших vision-language моделей (VLM) как более мощного семантического критерия. В работах по \textit{VLM-guidance} предлагается использовать функцию согласованности изображения и текста, вычисляемую через VLM, как сигнал для управления генерацией.

В частности, в работе Noise Diffusion вводится оценка $s(z_T)$, вычисляемая по изображению, полученному из начального шумового латента $z_T$ через процесс денойзинга, и вопросу о соответствии сэмпла промпту (VQA-score): $s(z_T)=\mathrm{VQA}(\cdot)$ \cite{miao2024noisediffusion}. Ключевая идея состоит в том, чтобы улучшать стартовый латент $z_T$, \emph{сохраняя распределение} $z_T \sim \mathcal{N}(0,I)$: вместо градиентного шага по $z_T$ применяется обновление, аналогичное прямому диффузионному процессу,
\[
z_T' = \sqrt{1-\gamma}\, z_T + \sqrt{\gamma}\, \sigma,
\qquad \sigma \sim \mathcal{N}(0,I),
\]
где $\gamma$~--- шаг (step size). В отличие от фиксированного шага, в работе $\gamma$ выбирается адаптивно в зависимости от текущего семантического качества:
\[
\gamma = 1 - \sqrt{s(z_T)}.
\]
Таким образом, при низком значении $s(z_T)$ шаг обновления становится более агрессивным, тогда как при высоком значении $s(z_T)$ обновление ослабляется.

Основная трудность --- подобрать шум $\sigma$, который увеличит семантическую оценку $s(\cdot)$. Для этого авторы в каждой итерации сэмплируют набор кандидатов $\{\sigma_i\}_{i=1}^N$, вычисляют соответствующие приращения
\[
v_i = z_T' - z_T
= (\sqrt{1-\gamma}-1)\,z_T + \sqrt{\gamma}\,\sigma_i,
\qquad i=1,\dots,N,
\]
и выбирают тот шум, для которого направление шага наиболее согласовано с оценкой градиента $\nabla_{z_T}s(z_T)$:
\[
\sigma^\star
= \arg\max_{\sigma_i}
\frac{\nabla_{z_T}s(z_T) v_i}{\lVert v_i\rVert_2^2}.
\]
После чего выполняется обновление
\[
z_T \leftarrow \sqrt{1-\gamma}\,z_T + \sqrt{\gamma}\,\sigma^\star,
\]
Такой шаг оптимизации происходит несколько раз итеративно.

\paragraph{VLM-gradient как обучающий сигнал для редактирования без пар.}
Отдельно важно направление, в котором VLM используется не только для inference-time guidance, но и как дифференцируемый сигнал при обучении модели редактирования без парных данных (оригинального и отредактированного изображения). В NP-Edit предлагается обучать диффузионную модель, опираясь на градиент VLM-оценки хода генерации отредактированной версии изображения \cite{kumari2025learning}.

\paragraph{Редактирование изображений и инверсия.}
В задачах редактирования дополнительным требованием является сохранение структуры исходного изображения, что во многом определяется точностью инверсии.

В работе Null-text Inversion показано, что оптимизация безусловного текстового эмбеддинга в рамках \textit{classifier-free guidance} позволяет значительно улучшить реконструкцию изображения и повысить устойчивость редактирования \cite{mokady2022nulltext}.

\paragraph{Stage-aware управление и proxy-prompts.}
Альтернативное направление представлено работами, в которых управление семантикой осуществляется не через градиенты, а через декомпозицию текстового условия по стадиям денойзинга. В частности, в подходе stage-aware prompt decomposition предлагается использовать последовательность proxy-prompts, адаптированных к различным фазам генерации \cite{huberman2025contextually}. Такой подход основан на наблюдении, что процесс диффузионной генерации носит \emph{coarse-to-fine} характер: на ранних шагах денойзинга формируется глобальная структура сцены и крупные композиционные элементы, тогда как на поздних шагах уточняются локальные детали, текстуры и мелкие атрибуты объектов.

Авторы показывают, что сложные или контекстно-противоречивые промпты могут перегружать модель на ранних этапах, когда необходимо прежде всего определить общую геометрию сцены. Поэтому предлагается сначала задавать более общие прокси-описания для формирования структуры, а затем постепенно уточнять семантику на более поздних шагах генерации. Такая стадийная декомпозиция позволяет частично повысить семантическую устойчивость и снизить вероятность противоречивых интерпретаций запроса.

Тем не менее данный подход воздействует на модель через изменение дискретного текстового условия и не обеспечивает прямого контроля над латентной траекторией. Кроме того, даже прокси-промпты могут оставаться семантически сложными для диффузионной модели, особенно в случаях количественных или нетипичных композиционных ограничений.

\paragraph{Позиционирование предлагаемого подхода.}
Существующие работы демонстрируют четыре ключевые идеи: (1) возможность управления диффузией через градиент внешней VLM-оценки верности семантики \cite{dhariwal2021diffusion,ho2022classifierfree}; (2) эффективность VLM как источника семантического сигнала \cite{miao2024noisediffusion,kumari2025learning}; (3) корень проблемы в воссоздании верной семантики именно ранних латентов \cite{huberman2025contextually}; (4) применимость оптимизации не только латентных непрерывных параметров для генерации \cite{mokady2022nulltext}.

В отличие от работ, где оптимизируется только начальный латент $z_0$, в данной работе предлагается исследовать оптимизацию промежуточных латентов $z_t$, а также непрерывных управляющих параметров (коэффициента CFG, null-text эмбеддингов и др.). Такой подход позволяет воздействовать на траекторию генерации на протяжении всего процесса денойзинга и потенциально обеспечивает более устойчивое выполнение сложных семантических ограничений как в задачах генерации, так и в задачах редактирования изображений.
